% =========================================================================
%  Hunarmand: technical dossier
%  A pitch document for the 5-minute jury round.
%
%  Compiled with lualatex so that fontspec can pick up JetBrains Mono and
%  EB Garamond. The brand palette tracks the website's globals.css.
% =========================================================================

\documentclass[11pt,a4paper]{article}

\usepackage[margin=2.4cm]{geometry}
\usepackage[english]{babel}
\usepackage{xcolor}
\usepackage{fontspec}
\usepackage{ebgaramond}
\usepackage{titlesec}
\usepackage{enumitem}
\usepackage{tcolorbox}
\usepackage{microtype}
\usepackage[hidelinks]{hyperref}
\usepackage{parskip}

% ── Palette ────────────────────────────────────────────────────────────
\definecolor{parchment}{HTML}{FAF7F2}
\definecolor{walnut}{HTML}{1C1410}
\definecolor{faded}{HTML}{5A4A3A}
\definecolor{muted}{HTML}{9A8A7A}
\definecolor{rouge}{HTML}{8B1A1A}
\definecolor{rougehover}{HTML}{C4454A}
\definecolor{gold}{HTML}{C8975A}
\definecolor{edge}{HTML}{D4C8B8}

% ── Page ───────────────────────────────────────────────────────────────
\pagecolor{parchment}
\color{walnut}

% ── Fonts ──────────────────────────────────────────────────────────────
% JetBrains Mono lives on disk under /usr/share/fonts/truetype/jetbrains-mono.
% We point fontspec at the .ttf files directly with Extension=.ttf so
% bold / italic / bold-italic resolution is unambiguous.
\setmonofont{JetBrainsMono}[
  Path=/usr/share/fonts/truetype/jetbrains-mono/,
  Extension=.ttf,
  Scale=0.90,
  Color=walnut,
  UprightFont=*-Regular,
  BoldFont=*-Bold,
  ItalicFont=*-Italic,
  BoldItalicFont=*-BoldItalic,
]

% ── Hyperref ───────────────────────────────────────────────────────────
\hypersetup{
  colorlinks=true,
  urlcolor=rouge,
  linkcolor=walnut,
  pdftitle={Hunarmand: technical dossier},
  pdfauthor={Bitwise},
}

% ── Section styling ────────────────────────────────────────────────────
\titleformat{\section}
  {\normalfont\Large\bfseries\color{walnut}}
  {\textcolor{rouge}{\thesection}}{0.6em}{}
\titleformat{\subsection}
  {\normalfont\large\bfseries\color{walnut}}
  {\textcolor{faded}{\thesubsection}}{0.6em}{}

\titlespacing*{\section}{0pt}{1.4\baselineskip}{0.5\baselineskip}
\titlespacing*{\subsection}{0pt}{1.0\baselineskip}{0.3\baselineskip}

% ── A small "stage" callout box used through the document ─────────────
\newtcolorbox{stage}[1]{
  enhanced,
  breakable,
  colback=parchment,
  colframe=edge,
  fonttitle=\bfseries\color{walnut},
  title={#1},
  boxrule=0.4pt,
  arc=2pt,
  left=8pt,right=8pt,top=6pt,bottom=6pt,
  before skip=8pt,after skip=8pt,
}

% ── A two-column key/value used in the "what we use" passages ─────────
\newcommand{\stack}[2]{%
  \par\noindent\textbf{\color{walnut}#1}\hspace{0.5em}\textcolor{rouge}{$\bullet$}\hspace{0.5em}\textcolor{faded}{#2}\par
}

% (no margin gloss helper needed; rules are inline where used)

% =========================================================================
\begin{document}
\pagestyle{empty}

\begingroup
  \raggedright
  {\fontsize{9pt}{11pt}\selectfont\color{faded}\texttt{tacit knowledge os \quad / \quad for kashmiri heritage artisans}}
  \par\vspace{0.4em}
  \textcolor{rouge}{\rule{2.6cm}{0.6pt}}
  \par\vspace{0.6em}
  {\fontsize{36pt}{40pt}\selectfont\bfseries hunarmand}
  \par\vspace{0.3em}
  {\fontsize{14pt}{18pt}\selectfont\itshape\color{faded}
  We treat what the master never wrote down as captureable, queryable, ownable infrastructure.}
  \par\vspace{1em}
  {\fontsize{9pt}{11pt}\selectfont\color{muted}\texttt{bitwise \quad / \quad srinagar 2026 \quad / \quad technical dossier v0.1}}
\endgroup

\vspace{1.2em}
\hrule height 0.4pt depth 0pt
\vspace{1.2em}

% =========================================================================
\section*{What this document is}
\addcontentsline{toc}{section}{What this document is}

A short, plain-English explanation of what we built, what each piece of
the platform does, and what technology we chose at every stage. The
companion deck is the 5-minute pitch; this is the dossier a juror
should be able to skim in 10 minutes if they want to know exactly how
the system works under the hood.

The writing avoids jargon where it can. Where a technical term is
unavoidable, we explain it the first time it appears. A short Q\&A at
the end answers the questions we expect a jury to ask.

% =========================================================================
\section{The system in one paragraph}

Hunarmand is three small services that share one Postgres database.
A \textbf{frontend} (Next.js, hosted on Vercel) for everything a user
sees. A \textbf{backend} (FastAPI on Render) that handles authentication,
masters, workshops, bookings, bundles, and orders. An \textbf{AI core}
(FastAPI on Hugging Face Spaces) that handles the things only AI can do:
speech recognition, structured extraction from interviews, embeddings,
and cryptographic signing of provenance certificates. The frontend
never talks to the AI core directly; everything routes through the
backend, so authentication is always honoured and the AI core stays
behind a single audit boundary.

The shared database is Neon Postgres with the \texttt{pgvector}
extension installed, which lets a single Postgres instance double as a
vector index for retrieval-augmented questions later in the flow.

% =========================================================================
\section{Architecture, in a single picture}

\begin{stage}{Three services, one database}
\begin{verbatim}
   user
     |
     v
   frontend (next.js, vercel)            <-- hunarmand.sal.lol
     |  https + jwt
     v
   backend (fastapi, render)             <-- hunarmand-backend.onrender.com
     |  https + bearer token       \
     v                              \
   ai core (fastapi, hf spaces)      \   <-- the slow expensive work
     |                                \
     +------> bhashini, ai4bharat, groq, hf inference, openai whisper
     +------> openrouter (llm aggregator), anthropic
     +------> sentence-transformers (local, no api call)
     +------> ed25519 + nacl secretbox (in-process)
     |
     v
   neon postgres + pgvector              <-- one database, both services
\end{verbatim}
\end{stage}

The arrows are HTTPS calls. JSON Web Tokens (JWTs) issued by the backend
authenticate the frontend; the backend forwards a bearer token to the
AI core for every call so the AI core knows which master is asking.

% =========================================================================
\section{Technology at each stage}

This section walks through every meaningful step a request goes through,
from the moment a master sits down with their phone to the moment a
buyer scans a QR on a finished shawl. Each stage explains \emph{what}
the step does, \emph{what we used} to do it, and \emph{why} we picked
that piece of technology.

% ─────────────────────────────────────────────────────────────────────
\subsection{Onboarding and authentication}

\begin{stage}{What happens}
A user types their phone number into the login screen. The frontend
asks the backend to send an OTP (one-time password). For the demo, the
OTP is fixed at \texttt{123456} so jurors can log in without an SMS
gateway in the loop. The user submits the OTP. The backend issues a
JWT, the frontend stores it in \texttt{localStorage}, and from that
point onward every request the user makes carries
\texttt{Authorization: Bearer <jwt>}.
\end{stage}

\stack{Auth}{phone-OTP login, JWT bearer tokens, role inferred from
phone (the demo artisan is \texttt{+919999999999}, anyone else is a
buyer)}
\stack{Why}{most heritage artisans have a phone but not a stable email.
WhatsApp and phone OTP are the de facto identity layer in Kashmir, so
we matched what already works.}

% ─────────────────────────────────────────────────────────────────────
\subsection{Vault: capturing what the master never wrote down}

This is the core innovation. The artisan sits in their own workshop and
the AI conducts a structured interview in Koshur (Kashmiri). The
interview is multi-pass: a broad pass to map the territory, then deeper
passes that drill into techniques, materials, lineage, environmental
adjustments, decision rules, and failure modes. The output is a
structured Craft DNA document that captures, in a queryable form, the
knowledge that would otherwise die with the master.

\subsubsection*{Stage 1: audio capture}

\begin{stage}{What happens}
The master speaks. The audio is uploaded to the AI core in chunks
(short windows), so a long interview can be transcribed and corrected
incrementally rather than waiting for a single huge file.
\end{stage}

\subsubsection*{Stage 2: speech-to-text (ASR)}

This is the trickiest stage because most speech-to-text engines do not
speak Koshur. We solved this with a \emph{ladder} of providers: each
engine is tried in order, and the first one that succeeds wins.

\stack{Bhashini}{the Government of India's national-language API. Has
first-class Kashmiri (\texttt{ks}) support. Free tier covers our scale
comfortably. \emph{This is our primary ASR.}}
\stack{AI4Bharat IndicWhisper}{open-weight Indic ASR from IIT Madras.
Used as a backup if Bhashini is unavailable. Self-hostable, so we can
move it in-house later.}
\stack{Groq Whisper}{a hosted Whisper variant on Groq's free tier.
Catches Koshur surprisingly well via Urdu acoustic similarity, and
returns word-level timestamps.}
\stack{Hugging Face Inference Whisper}{another fallback path; cold
starts are slow but it costs nothing.}
\stack{OpenAI Whisper}{the paid premium fallback if every free option
is down.}
\stack{Manual transcript}{if all engines fail, an operator can submit
a typed transcript. Nothing in the rest of the pipeline depends on a
specific provider.}

\stack{Why a ladder}{the demo costs zero rupees a month, but if traffic
spikes or one provider rate-limits us, the ladder absorbs the spike
without operator intervention.}

\subsubsection*{Stage 3: translation, when needed}

If the transcript is in Koshur or Urdu, an LLM translates it to English
before structured extraction. We use the same provider abstraction as
the LLM (next stage), so the translation step is just a different
prompt against the same client.

\stack{Translation}{routed through the LLM client; Koshur to English is
the common path.}

\subsubsection*{Stage 4: structured extraction (Craft DNA)}

We do \emph{not} ask an LLM to "summarize" the interview. Instead, we
constrain the LLM with a strict JSON schema we call Craft DNA. The
schema has fields like \texttt{technique\_names},
\texttt{materials\_summary}, \texttt{lineage\_chain},
\texttt{environmental\_tunings}, \texttt{decision\_rules}, and
\texttt{failure\_log}. Each pass focuses on filling one slice of the
schema with citations back into the transcript.

\stack{LLM provider}{OpenRouter, which aggregates frontier-class models
(including free-tier ones) behind one OpenAI-compatible interface. We
swap models with a single environment variable.}
\stack{Adaptive structured output}{some models support
\texttt{json\_schema} response format natively, others only
\texttt{json\_object}, others nothing. Our generic structured-output
helper tries each strategy in order and falls back automatically.}
\stack{Repair loop}{if the model returns malformed JSON or violates the
schema, we feed the validator's error message back to the model and
ask it to fix the specific problem. Continues for a fixed number of
retries before giving up.}
\stack{Why schema-first}{LLMs are unreliable when asked to write free
text. They are remarkably reliable when asked to fill a slot. Forcing
the model to fill a Pydantic-validated schema turns the AI from a
storyteller into a typewriter.}

\subsubsection*{Stage 5: embeddings}

Every transcript chunk is embedded. We use \texttt{sentence-transformers},
specifically the \texttt{intfloat\allowbreak/multilingual\allowbreak-e5\allowbreak-small}
model, which produces 384-dimensional vectors and runs locally in the
AI core container. No external API call. No per-vector cost.

\stack{Model}{intfloat/multilingual-e5-small via sentence-transformers}
\stack{Where}{runs locally inside the AI core (downloaded once at
container build time, cached forever after)}
\stack{Why}{free, fast on CPU, multilingual (covers Koshur reasonably),
no rate limits, no surprise bills}

\subsubsection*{Stage 6: storage and indexing}

\begin{stage}{What lives where}
The Craft DNA document goes into the \texttt{craft\_dnas} table. Each
transcript chunk goes into \texttt{vault\_chunks} with its embedding,
its timestamps, and a foreign key back to the master. The
\texttt{embedding} column is a \texttt{pgvector} column, so we can run
cosine-similarity searches in SQL with no separate vector store.
\end{stage}

\stack{Database}{Neon Postgres with pgvector. One database for both
the backend and the AI core. We rename the AI core's tables to
\texttt{ai\_*} where they would otherwise collide with backend tables.}
\stack{Migrations}{Alembic, run automatically on every backend deploy
via a \texttt{start.sh} entrypoint that calls \texttt{alembic upgrade
head} before \texttt{uvicorn}.}

% ─────────────────────────────────────────────────────────────────────
\subsection{Sanad: cryptographic provenance}

A Sanad is a verifiable certificate that says "this specific physical
piece was made by this specific master's lineage". When a master
finishes a piece, they mint a Sanad for it from the studio. The result
is a small QR code that gets attached to the piece.

\subsubsection*{The cryptographic stack, step by step}

\stack{Step 1, keys}{the AI core generates an Ed25519 keypair for each
master. The public key is published. The private key is encrypted at
rest using NaCl SecretBox with a key-encryption-key derived from a
process-level secret. Stored in \texttt{master\_keys}.}
\stack{Step 2, payload}{the master fills a small form (piece id, craft,
techniques, materials, lineage). The form data, plus issuance
timestamp, becomes the JSON payload of the Sanad.}
\stack{Step 3, canonical bytes}{the JSON payload is run through RFC
8785 \emph{JSON Canonicalization Scheme} (JCS). JCS guarantees that
two semantically identical JSON objects produce byte-identical output,
regardless of key order or whitespace. This is critical: signatures
are over bytes, so we need a deterministic byte form.}
\stack{Step 4, signature}{the canonical bytes are signed with the
master's Ed25519 private key. Output is a 64-byte signature.}
\stack{Step 5, JWS-compact}{header.payload.signature, each component
base64url-encoded, joined with dots. This is the format a buyer's
verifier expects, and it is a single short string.}
\stack{Step 6, QR}{the JWS-compact string is rendered as a PNG QR code
and embedded in the response. The frontend renders it; the master
prints it.}

\subsubsection*{Verification, offline}

A buyer scans the QR. Their verifier app gets the JWS string, splits
it on dots, base64-decodes each component, runs Ed25519 verify against
the master's public key (cached or fetched from the public registry).
If the signature checks, the piece is authentic. If any byte of the
payload is altered, verification fails.

\stack{Why offline}{signatures are mathematics, not platform trust. A
verifier with the public key cached can validate a Sanad on a flight,
in a remote village, or after Hunarmand goes out of business in 50
years. The signature outlives the platform.}
\stack{Forgery resistance}{a counterfeit shawl can be made. A Sanad
signed by Mohammad Yusuf cannot be forged without his private key,
which lives encrypted in our service and is never released.}

% ─────────────────────────────────────────────────────────────────────
\subsection{Ustaad: workshops}

Standard CRUD over a \texttt{workshops} table. An artisan logged in
through the studio can create, edit, hide, or delete their offerings.
A buyer browses the public listing, picks a slot, and books with their
phone number; the booking lands in the \texttt{bookings} table.

\stack{Stack}{FastAPI handlers, async SQLAlchemy ORM, Pydantic schemas,
Postgres as the source of truth. No magic.}

% ─────────────────────────────────────────────────────────────────────
\subsection{Bazaar: commerce}

An artisan curates Bundles (sets of pieces, optionally linked to
Sanads). A buyer adds bundles to their cart and checks out with their
phone number. For the demo we wire a no-op checkout that creates the
\texttt{Order} row and shows a confirmation page; production would
attach Stripe (the SDK is already pulled in) for actual payment.

\stack{Stack}{same as Ustaad, plus a thin checkout handler. Stripe SDK
is configured but optional.}

% ─────────────────────────────────────────────────────────────────────
\subsection{Ask the Hunarmand: retrieval-augmented questions}

This is how the master's tacit knowledge becomes queryable by anyone.
A user types a question; the platform returns an answer with citations
back into the master's own words.

\subsubsection*{Step by step}

\stack{1.\ embed the question}{the question runs through the same
sentence-transformers model used at ingest time, producing a
384-dimensional query vector.}
\stack{2.\ retrieve}{a SQL query against \texttt{vault\_chunks} ordered
by cosine distance returns the top-K closest chunks for the master in
question. \texttt{pgvector}'s SQLAlchemy adapter binds the vector
parameter as a real \texttt{Vector} type, not a string.}
\stack{3.\ generate, with strict guardrails}{the chunks plus the
question are passed to an LLM with explicit instructions: only answer
using the chunks, cite each claim with the chunk's start timestamp,
and refuse if the chunks do not contain a clear answer. Refusal is a
first-class response, not a hallucination.}
\stack{4.\ return}{the answer object contains the prose answer, the
citations (chunk id + timestamp), and a confidence score.}

\stack{Why the strict refusal}{the failure mode for a craft-knowledge
system is a confident wrong answer about a master's technique. We log
refusals as a quality signal, not as a bug.}

% =========================================================================
\section{Why these choices, in plain English}

\begin{itemize}[leftmargin=*]
  \item \textbf{Free-tier first.} Bhashini, AI4Bharat, OpenRouter free
        models, sentence-transformers on CPU, Neon Postgres, Vercel,
        Render, Hugging Face Spaces. The entire demo runs at zero
        rupees a month. Every paid path is a swap of one environment
        variable.
  \item \textbf{Provider abstraction.} ASR, LLM, embeddings, and storage
        each sit behind a uniform interface. If Groq goes down, Whisper
        picks up. If OpenRouter rate-limits us, Anthropic fills in.
        Nothing in the application layer cares which provider answered.
  \item \textbf{Schema-first AI.} We never let an LLM write free text
        for the Craft DNA. Every output is constrained by a strict
        JSON schema with a repair loop. The model becomes a slot-filler,
        which it is reliably good at.
  \item \textbf{Offline verification.} Sanads use Ed25519 signatures
        over RFC 8785 canonical JSON, encoded as JWS-compact strings,
        rendered as QR codes. A verifier with the master's public key
        can validate a Sanad without ever talking to our servers.
  \item \textbf{Indigenous tradition.} The Craft DNA schema is the
        same shape as the Kashmiri Talim, an encoded weaving notation
        the master kept on paper centuries before code existed. We are
        not bringing foreign technology to Kashmir. We are upgrading
        what Kashmir already invented.
\end{itemize}

% =========================================================================
\section{Anticipated questions from the jury}

A short Q\&A on questions we expect to land. Each answer is honest
about scope: this is a hackathon build, not a finished business.

\begin{description}[style=nextline,leftmargin=0pt,itemsep=0.5em]

\item[Why Postgres + pgvector instead of a dedicated vector database?]
One database is easier to operate, cheaper, and pgvector is plenty
fast for our scale. Pinecone or Weaviate would mean another service,
another bill, another point of failure. Neon's free tier covers our
entire vector index today and well beyond.

\item[How do you stop the AI from hallucinating Craft DNA?]
Two mechanisms. First, every LLM call is constrained by a strict JSON
schema; the model cannot return free text. Second, the Ask-the-Hunarmand
endpoint refuses to answer questions where the retrieved chunks do not
contain a clear answer. Refusals are logged as a quality signal.

\item[What about Koshur? Most LLMs do not speak it.]
We use Bhashini and AI4Bharat IndicWhisper for Koshur ASR. Both have
first-class Kashmiri support. For LLM steps, we route Koshur through
translation to English before structured extraction; the structured
output is then sent back through the buyer-facing flow in English with
the original Koshur preserved as a citation.

\item[Can a counterfeiter sign a fake Sanad in a master's name?]
No. The Sanad is signed by the master's private key, which is stored
encrypted at rest on our AI core and never released. The public key is
published. A counterfeiter can copy the QR off a real piece, but the
signature only validates against the original payload bytes. Any
change (different master, different piece, altered date) breaks the
signature.

\item[Why phone OTP instead of email?]
Most artisans we'd onboard do not have a stable email. They all have a
phone. WhatsApp is the de facto identity layer in Kashmir, so we
matched what already works.

\item[How do you prevent a master from over-signing, flooding the
market with their own forged pieces?]
Technically, we cannot. But the platform records every Sanad ever
signed by every master. A buyer can see how many pieces a master has
signed, how often, and the full provenance chain. Reputation is on
the public registry. A master who over-signs degrades their own market
in a way the registry makes obvious.

\item[What is the business model?]
A platform commission of 6 to 9 percent on direct commerce, against
the 70 to 85 percent that the existing middleman chain currently
extracts. The artisan keeps the rest.

\item[How do you scale to ten thousand masters?]
The ASR ladder absorbs traffic spikes (free providers swap to paid as
needed). Postgres + pgvector scales vertically into the millions of
vectors. The signing service is stateless and horizontally scalable.
The bottleneck is likely the field team onboarding masters, not
infrastructure.

\item[What if Bhashini changes their API?]
Every provider in the ladder is abstracted behind a uniform interface.
If Bhashini changes, we swap that single implementation. The rest of
the system does not notice.

\item[How long does an interview take?]
About two to three hours per master, conducted across two or three
sessions. The structured AI interview keeps the master on track and
surfaces details they would not volunteer in a free-form conversation.

\item[Have you talked to actual masters?]
Hackathon scope means we built the platform, not the field operation.
The vision is to deploy with NGO partners in the valley for actual
onboarding once the platform is stress-tested.

\item[Why not just use a Google Form for the interview?]
Because the master does not read English. Because the master cannot
write. Because the master answers questions verbally, in Koshur, the
way they would teach an apprentice. The AI interview is the first
time their knowledge can be captured \emph{as they speak it}.

\item[How do disputes work, what if a master's family says a Sanad
should not have been signed?]
Each master controls their own keypair. A new keypair version
supersedes an old one for new pieces; old Sanads continue to verify
against their original key (looked up by \texttt{kid}). The platform
is a ledger, not an authority.

\item[What stops the platform itself from being the new middleman?]
The platform commission cap of 6 to 9 percent is structural, not
discretionary. The Sanad is offline-verifiable, so the signature
does not depend on us being alive. The Craft DNA is owned by the
master, who can export the structured artefact and take it elsewhere.

\item[How long until you can launch?]
The platform is functional today. Field deployment depends on
partnerships with valley-based NGOs and on getting Bhashini's
production API limits raised.

\item[What is the single most ambitious thing in this build?]
The structured interview that turns a craft master's spoken answers
into a 240-field schema, signed under their own keypair, indexed for
retrieval, owned by the master in perpetuity. That artefact did not
exist before this hackathon. It does now.

\end{description}

\vspace{1em}
\hrule height 0.4pt depth 0pt
\vspace{0.6em}

\begingroup
  \centering
  \fontsize{10pt}{12pt}\selectfont\color{faded}
  \texttt{hunarmand \quad / \quad bitwise \quad / \quad srinagar 2026}\par
\endgroup

\end{document}
